\section*{Chapitre \ref{ch:qm:bayesian-methods} --- Méthodes bayésiennes}

\begin{solution}{exo:qm:bayesian-methods:1}
$\mathrm{Beta}(80, 70)$, moyenne $80/150 = 0.533$. Avec la \omterm{def:qm:bayesian-methods:prior}{loi a priori} plate, $\mathrm{Beta}(31, 21)$, moyenne $31/52 = 0.596$~: la \omterm{def:qm:bayesian-methods:prior}{loi a priori} informative
ramène presque entièrement le taux de réussite de 60~\% vers celui d'une pièce équilibrée.
\end{solution}

\begin{solution}{exo:qm:bayesian-methods:2}
$\sigma^2 = 0.2$, $\tau^2 = 0.16$, $B = 0.2/0.36 = 0.556$. Moyenne a posteriori $0.5 + 0.444 \times 1.0 = 0.94$~; variance a posteriori $0.2 \times
0.16/0.36 = 0.089$, écart-type 0.30.
\end{solution}

\begin{solution}{exo:qm:bayesian-methods:3}
$n(1 - \rho)/(1 + \rho) = 10\,000 \times 0.2/1.8 = 1\,111$.
\end{solution}

\begin{solution}{exo:qm:bayesian-methods:4}
$B = \sigma^2/(\sigma^2 + \hat\tau^2) = \sigma^2/s_x^2$, donc l'estimation conserve $1 - \sigma^2/s_x^2$ de chaque écart. L'\omterm{def:qm:bayesian-methods:js}{estimateur de James--Stein} vers la
moyenne globale conserve $1 - (k - 3)\sigma^2/\sum(x_i - \bar x)^2 = 1 - \frac{k-3}{k-1}\sigma^2/s_x^2$~: le même \omterm{def:qm:estimation:estimator}{estimateur} avec un rétrécissement
un peu moindre, l'écart s'annulant à mesure que $k$ croît.
\end{solution}

\begin{solution}{exo:qm:bayesian-methods:5}
$B = \sigma^2/(\sigma^2 + \tau^2)$ avec $\sigma^2 = 1/T$ vaut la moitié lorsque $1/T = \tau^2$~: avec $\hat\tau^2 = 0.168$, $T = 6.0$ ans (6.25 avec la
vraie $\tau = 0.4$). Il faut six ans d'historique avant que celui d'un gérant ne pèse autant que la \omterm{def:qm:bayesian-methods:prior}{loi a priori} de la plateforme.
\end{solution}

\begin{solution}{exo:qm:bayesian-methods:6}
Sachant $\theta$ et $m$, la vraisemblance de $\tau^2$ vaut $\prod_i\tau^{-1}e^{-(\theta_i - m)^2/(2\tau^2)} = (\tau^2)^{-k/2}e^{-S/(2\tau^2)}$ avec $S =
\sum_i(\theta_i - m)^2$. Une \omterm{def:qm:bayesian-methods:prior}{loi a priori} plate sur $\tau$ est $p(\tau^2) \propto (\tau^2)^{-1/2}$, donc la loi conditionnelle est proportionnelle à
$(\tau^2)^{-(k-1)/2 - 1}e^{-S/(2\tau^2)}$~: une loi gamma inverse de forme $(k - 1)/2$ et d'échelle $S/2$.
\end{solution}

\begin{solution}{exo:qm:bayesian-methods:7}
\texttt{best\_manager\_average()}~: le meilleur historique affiche en moyenne 2.09, le \omterm{def:qm:estimation:sharpe}{ratio de Sharpe} réel du gérant retenu vaut 1.01, l'estimation par \omterm{def:qm:bayesian-methods:hier}{bayésien empirique}
vaut 1.01 et les trois années suivantes 1.02. La sélection ajoute 1.08 à l'historique~; l'estimation rétrécie la retire en moyenne.
\end{solution}

\begin{solution}{exo:qm:bayesian-methods:8}
Rétrécir vers zéro avec le facteur de James--Stein revient à traiter les signaux comme des tirages issus d'une même population centrée en zéro, et réduit bien
l'erreur moyenne des estimations. Elle n'annule pas la recherche~: le facteur dépend de la dispersion des \omterm{def:qm:estimation:sharpe}{ratios de Sharpe} publiés, et si seuls les
survivants d'une recherche plus large sont publiés, cette dispersion et ce point de référence sont tous deux faux, si bien que les signaux retenus restent surestimés. La \omterm{def:qm:bayesian-methods:prior}{loi a priori}
doit être calibrée sur l'ensemble des essais, et les tests du chapitre~12 restent applicables.
\end{solution}

\begin{solution}{pb:qm:bayesian-methods:1}
\textbf{1.} Meilleur 2.41~; moyenne 0.51, écart-type 0.71.
\textbf{2.} Douze.
\textbf{3.} $1/\sqrt3 = 0.58$~: les historiques se dispersent davantage que le seul bruit ne l'explique.
\textbf{4.} $\hat m = 0.51$, $\hat\tau = \sqrt{0.71^2 - 1/3} = 0.41$.
\textbf{5.} $B = 0.333/(0.333 + 0.168) = 0.665$.
\textbf{6.} $0.51 + 0.335 \times 1.90 = 1.14$, écart-type 0.33.
\textbf{7.} 1.20 (facteur 0.36 au lieu de 0.335).
\textbf{8.} 1.82 et 1.45 deviennent 0.95 et 0.83~; le classement est inchangé, puisque le rétrécissement est le même pour tous.
\textbf{9.} Ratios réels 1.07, 0.96 et 1.68~; sur les trois années suivantes, 1.69, 0.57 et 0.99. Le meilleur historique n'était pas celui du meilleur gérant.
\textbf{10.} Loi prédictive $\mathcal N(1.14, 0.67^2)$~: $\Phi(-1.14/0.67) = 4.3\%$.
\textbf{11.} De 0.335 à 0.149, soit une baisse de 55~\%.
\textbf{12.} Par rapport au ratio réel, de 0.336 à 0.121 (64~\%)~; par rapport aux trois années suivantes, de 0.672 à 0.457 (32~\%).
\textbf{13.} Meilleur historique 2.09, ratio réel 1.01, estimation rétrécie 1.01 (trois années suivantes 1.02).
\textbf{14.} Moyenne a posteriori 1.12, \omterm{def:qm:bayesian-methods:predictive}{intervalle de crédibilité} à 95~\% $(0.41, 2.01)$~: plus large que l'intervalle de substitution, car $\tau$ est incertain.
\textbf{15.} Moyenne 0.40, intervalle à 95~\% $(0.13, 0.66)$, avec une probabilité de 7~\% que $\tau < 0.2$, auquel cas les gérants se distinguent à peine et le
rétrécissement devrait être bien plus fort.
\textbf{16.} Pas sur la seule base de ces données~: le \omterm{def:qm:estimation:sharpe}{ratio de Sharpe} attendu vaut 1.14, très bon mais pas 2.4, et le capital doit être dimensionné sur 1.14
en tenant compte de son incertitude, avec des apports supplémentaires à mesure que l'historique s'allonge.
\textbf{17.} Une \omterm{def:qm:bayesian-methods:prior}{loi a priori} calibrée sur des gérants externes comparables, plus large (avec un $\tau$ plus grand) et probablement centrée plus bas, puisque le recrutement se fonde
aussi sur les historiques passés.
\textbf{18.} Lorsque les \omterm{def:qm:estimation:sharpe}{ratios de Sharpe} réels ont des queues lourdes, quelques gérants réellement exceptionnels existent et la \omterm{def:qm:bayesian-methods:prior}{loi a priori} normale les ramène aussi fortement
que les gérants chanceux~; une \omterm{def:qm:bayesian-methods:prior}{loi a priori} de Student-$t$ rétrécit moins les grands écarts.
\textbf{19.} Résultat nommé~: \emph{la présélection de l'allocateur}~: le meilleur parmi cinquante historiques de trois ans, 2.41, a une moyenne a posteriori
de 1.14 selon le \omterm{def:qm:bayesian-methods:hier}{bayésien empirique}~; le rétrécissement réduit de 31~\% l'erreur par rapport aux historiques des trois années suivantes sur cette plateforme, et de 32~\% en moyenne.
\textbf{20.} Parce que le meilleur parmi de nombreux historiques bruités est sélectionné en partie pour son bruit, et que la \omterm{def:qm:bayesian-methods:prior}{loi a priori} retire la part de l'écart que
le bruit explique.
\end{solution}

\begin{solution}{iq:qm:bayesian-methods:1}
Un intervalle de confiance à 95~\% est produit par une procédure qui couvre le paramètre vrai, fixe, dans 95~\% des échantillons répétés~; un \omterm{def:qm:bayesian-methods:predictive}{intervalle de crédibilité} à 95~\%
contient le paramètre avec une probabilité a posteriori de 0.95 sachant cet échantillon et la \omterm{def:qm:bayesian-methods:prior}{loi a priori}. Ils coïncident numériquement dans les modèles simples à \omterm{def:qm:bayesian-methods:prior}{loi a priori} plate
mais diffèrent toujours par leur sens.
\iqlookfor{Fréquence sur les échantillons opposée à probabilité sachant les données, et rôle de la \omterm{def:qm:bayesian-methods:prior}{loi a priori}.}
\end{solution}

\begin{solution}{iq:qm:bayesian-methods:2}
Traiter les cinquante gérants comme des tirages d'une population~: estimer sa moyenne et sa dispersion à partir des historiques (en retranchant la variance du bruit), et rétrécir
chaque historique vers la moyenne par $\sigma^2/(\sigma^2 + \tau^2)$. Avec trois ans et une dispersion typique, les deux tiers de l'excès du meilleur historique
sont du bruit~; il faut s'attendre à environ la moitié de son \omterm{def:qm:estimation:sharpe}{ratio de Sharpe}.
\iqlookfor{Régression vers la moyenne quantifiée, variance du bruit $1/T$, et \omterm{def:qm:bayesian-methods:hier}{bayésien empirique}.}
\end{solution}

\begin{solution}{iq:qm:bayesian-methods:3}
Parce que l'erreur quadratique totale sur trois moyennes ou plus peut être réduite en échangeant un peu de biais contre une forte baisse de la variance~; tirer chaque
estimation vers un point commun y parvient quelles que soient les vraies moyennes. Le piège~: le gain porte sur le total et non sur chaque unité, et une
unité inhabituelle peut en pâtir.
\iqlookfor{Compromis biais-variance en dimension trois ou plus, et cette domination porte sur la somme.}
\end{solution}

\begin{solution}{iq:qm:bayesian-methods:4}
On ne le sait jamais~; on le vérifie. Lancer plusieurs chaînes à partir de points de départ dispersés et les comparer ($\hat R$ proche de 1), examiner les tracés, écarter une
\omterm{def:qm:bayesian-methods:ess}{période de chauffe} (burn-in), calculer les \omterm{def:qm:bayesian-methods:ess}{tailles d'échantillon effectives} de chaque quantité rapportée, et comparer avec un autre échantillonneur ou un cas particulier analytique.
\iqlookfor{Plusieurs chaînes, $\hat R$, \omterm{def:qm:bayesian-methods:ess}{taille d'échantillon effective}, et prudence quant à la convergence.}
\end{solution}

\begin{solution}{iq:qm:bayesian-methods:5}
L'\omterm{def:qm:estimation:estimator}{estimateur} de la régression ridge est le mode a posteriori (et la moyenne) d'un modèle linéaire muni d'une \omterm{def:qm:bayesian-methods:prior}{loi a priori} normale $\beta \sim \mathcal N(0, (\sigma^2/\lambda)I)$ sur
les coefficients~: la pénalité $\lambda\lvert\beta\rvert^2$ est l'opposé du logarithme de la \omterm{def:qm:bayesian-methods:prior}{loi a priori}. Le lasso est le mode sous une \omterm{def:qm:bayesian-methods:prior}{loi a priori} de Laplace.
\iqlookfor{Pénalité comme opposé du logarithme de la \omterm{def:qm:bayesian-methods:prior}{loi a priori}, et échelle de $\lambda$ comme variance a priori.}
\end{solution}

\begin{solution}{iq:qm:bayesian-methods:6}
Sa règle d'acceptation égalise le flux de probabilité de $x$ vers $y$ et le flux de $y$ vers $x$ sous $\pi$ (\omterm{def:qm:markov-chains-and-queues:balance}{équilibre détaillé}), si bien que $\pi$ est une
\omterm{def:qm:stochastic-differential-equations:kolmogorov}{loi stationnaire}, et une chaîne ergodique converge vers elle. Un pas de Gibbs propose à partir de la loi conditionnelle exacte, ce qui rend le rapport de Metropolis--Hastings
exactement égal à un.
\iqlookfor{\omterm{def:qm:markov-chains-and-queues:balance}{Équilibre détaillé} explicité, et simplification dans le cas de Gibbs.}
\end{solution}
